\section{高级调试}

高级调试区域默认是收起的，点"高级调试"按钮才会展开。\textbf{普通售后调试不建议使用，除非研发明确要求。}

\begin{figure}[H]
    \centering
    \includegraphics[width=0.8\textwidth]{figures/ui-advanced.png}
    \caption{高级调试区域（展开状态）。包含 PWM 和速度校准 RPM。}
    \label{fig:ui-advanced}
\end{figure}

\subsection{PWM}

PWM（脉宽调制）是底层速度控制值。它和界面上的"速度"是同一回事的两种表达方式：

\begin{itemize}
    \item "速度"更直观，售后日常用这个。
    \item "PWM"更底层，数值范围和物理含义不那么直观。
\end{itemize}

\textbf{注意：本固件里 PWM 数值越大，速度越慢；PWM 数值越小，速度越快。}这和直觉相反，容易搞混，所以售后优先调"速度"，不要直接调 PWM。

\begin{itemize}
    \item \textbf{默认值：}\texttt{150}
    \item \textbf{可调范围：}\texttt{0 \textasciitilde{} 500}
\end{itemize}

\begin{table}[H]
    \centering
    \caption{PWM 按钮说明}
    \label{tab:pwm-buttons}
    \begin{tabular}{lp{6cm}}
        \toprule
        \textbf{按钮} & \textbf{作用} \\
        \midrule
        应用 & 临时设置 PWM，不保存 \\
        保存 & 保存 PWM，断电后保留 \\
        \bottomrule
    \end{tabular}
\end{table}

\subsection{速度校准 RPM}

\begin{itemize}
    \item \textbf{默认值：}\texttt{3655 RPM}
    \item \textbf{作用：}用于"速度 cm/min"和底层 PWM 之间的换算
\end{itemize}

一般不要修改。只有研发确认电机型号、轮径、传动结构或测速基准发生变更后，才需要改这个值。
